

\section{Lògica i raonament}
\iniciTest{t1}{c,d,b,c,b,a,d,b,b,a}

\begin{enumerate}
\pregunta Considerem l'enunciat següent sobre una festa: "Si és el teu
aniversari o hi haurà pastís, hi haurà pastís". Quina de les
següents afirmacions és falsa?

\begin{enumerate}
\opcio{a}{ Si $p=$\enquote{és el teu aniversari}\ i
$q=$\enquote{ hi haurà pastís}, llavors
l'enunciat s'expressa formalment així: $\left(  p\vee q\right)
\longrightarrow q$.

}
\opcio{b}{ Suposant que l'enunciat és cert, no podem concloure necessàriament que hi haurà pastís.

}
\opcio{c}{ Suposant que l'enunciat és cert, si no hi haurà pastís
podem concloure que és el teu aniversari. 

}
\opcio{d}{ Suposem que s'ha descobert que l'afirmació és mentida, aleshores
podem concloure que és el teu aniversari, però el pastís és mentida.
}
\end{enumerate}
\feedback

\pregunta Si $p=$\enquote{ Joan sempre diu la veritat}\ i
$q=$\enquote{ Jaume sempre menteix}, llavors quina de
les següents afirmacions és falsa?

\begin{enumerate}
\opcio{a}{ L'enunciat \enquote{ No és cert que en Joan o Jaume diguin
sempre mentides}\ s'expressa formalment d'aquesta manera:
$\lnot\left(  \lnot p\vee q\right)  $.

}
\opcio{b}{ $\lnot p\longrightarrow\lnot q$ és l'enunciat \enquote{ Si
en Joan menteix, llavors en Jaume diu la veritat}.

}
\opcio{c}{ $p\wedge\lnot q$ és l'enunciat \enquote{ Joan i Jaume no
menteixen}.

}
\opcio{d}{ L'enunciat \enquote{ No és el cas que si en Joan menteix,
aleshores en Jaume digui sempre la veritat}\ s'expressa
formalment així $\lnot\left(  \lnot p\longrightarrow q\right)  $. 
}
\end{enumerate}
\feedback

\pregunta Prenent com univers el conjunt dels nombres enters $\mathbb{Z}$. Si
$p=$\enquote{$x$ és múltiple de $3$}\ i
$q=$\enquote{$x$ és enter positiu més petit que $10$%
}. Quina de les següents afirmacions és vertadera?

\begin{enumerate}
\opcio{a}{ $\lnot\left(  \lnot p\vee\lnot q\right)  =$\enquote{$x$ és
múltiple de $3$ més petit que $10$}.

}
\opcio{b}{ $p\longrightarrow\lnot q=$\enquote{ si $x$ és múltiple de
$3$, aleshores $x$ és un enter més petit o igual que $0$ o és
més gran o igual que $10$}. 

}
\opcio{c}{ $\left(  p\wedge\lnot q\right)  \vee\left(  \lnot p\wedge q\right)
=$\enquote{$x$ és múltiple de $3$ o $x$ és enter positiu
més petit que $10$}.

}
\opcio{d}{ $\left(  p\longrightarrow q\right)  \vee\lnot p$ és una
proposició que només és falsa quan assignem a la variable $x$
valors enters més grans o iguals que $10$.
}
\end{enumerate}
\feedback

\pregunta La proposició $\left(  \lnot p\longleftrightarrow\lnot q\right)
\vee\left(  r\longrightarrow q\right)  $ és falsa quan:

\begin{enumerate}
\opcio{a}{ $p$ i $q$ són falses, i $r$ és vertadera.

}
\opcio{b}{ $p$ és vertadera, i $q$ i~$r$ són vertaderes.

}
\opcio{c}{ $p$ i $r$ són vertaderes, i $q$ és falsa. 

}
\opcio{d}{ $p,q$ i $r$ són falses.
}
\end{enumerate}
\feedback

\pregunta Quin parell de formes proposicionals són lògicament equivalents?

\begin{enumerate}
\opcio{a}{ $A\wedge B$ i $\lnot\left(  \lnot A\longrightarrow B\right)  $

}
\opcio{b}{ $\left(  A\wedge B\right)  \vee\left(  B\vee\lnot C\right)  $ i $\left(
\lnot A\vee\lnot B\right)  \longrightarrow\lnot\left(  \lnot B\wedge C\right)
$ 

}
\opcio{c}{ $\lnot\left(  (A\rightarrow\lnot B)\vee\lnot(C\wedge\lnot C)\right)  $ i
$C\vee\lnot C$

}
\opcio{d}{ $\left(  A\longrightarrow B\right)  \longleftrightarrow B$ i $A$
}
\end{enumerate}
\feedback

\pregunta En Pere t'estava explicant el que va menjar ahir a la tarda. Et diu:
\enquote{ Vaig prendre crispetes o panses. A més, si tenia
entrepans de cogombre, llavors tenia refresc. Però no vaig beure refresc
ni te}. Per descomptat, sabeu que en Pere és el pitjor mentider del
món, i tot el que diu és fals. Què va menjar en Pere?

\begin{enumerate}
\opcio{a}{ entrepans de cogombre\thinspace\ i te 

}
\opcio{b}{ crispetes o panses i refresc

}
\opcio{c}{ només te

}
\opcio{d}{ entrepans de cogombre\thinspace\ i refresc
}
\end{enumerate}
\feedback

\pregunta De les premisses
\[
P_{1}=\lnot A\longrightarrow\left( B\longrightarrow\lnot C\right),\quad
P_{2}=C\longrightarrow\lnot A,\quad
P_{3}=\left( \lnot D\vee A\right) \longrightarrow\lnot\lnot C
\]
i $P_{4}=\lnot D$, què es pot deduïr si és vàlid?

\begin{enumerate}
\opcio{a}{ Podem deduir qualsevol cosa perquè les premisses són inconsistents.

}
\opcio{b}{ $B\wedge\lnot\left(  A\vee D\right)  $

}
\opcio{c}{ $\lnot B\wedge\left(  A\vee D\right)  $

}
\opcio{d}{ $\lnot\left(  A\vee B\vee D\right)  $ 
}
\end{enumerate}
\feedback

\pregunta Volem simplificar els enunciats següents de manera que la
negació només apareix directament al costat dels predicats, quina de les
següents equivalències no és correcta?

\begin{enumerate}
\opcio{a}{ $\lnot\exists x\forall y(\lnot O(x)\vee E(y))\Longleftrightarrow\forall
x\exists y(O(x)\wedge\lnot E(y))$

}
\opcio{b}{ $\lnot\forall x\lnot\forall y\lnot(x<y\wedge\exists z(x<z\vee
y<z))\Longleftrightarrow\exists x\forall y(x<y\wedge\exists z(x<z\vee y<z)) $ 

}
\opcio{c}{ \enquote{ Hi ha un nombre $n$ per al qual cap altre nombre
és menor o igual a $n$}\ és equivalent a
\enquote{ Hi ha un nombre $n$ per al qual tots els altres nombres
són estrictament majors que $n$}.

}
\opcio{d}{ \enquote{És fals que per a cada nombre $n$ hi hagi dos
nombres més entre els quals es troba $n$}\ és
equivalent a \enquote{Hi ha un nombre $n$ que no està entre
altres dos nombres}.
}
\end{enumerate}
\feedback

\pregunta Es diu que el raonament següent és vàlid%
\[%
\begin{tabular}
[c]{ll}%
$P_{1}:$ & $\exists x~\left(  P(x)\wedge Q(x)\right)  $\\
$P_{2}:$ & $\exists x$ $\left(  M(x)\wedge\lnot P(x)\right)  $\\\hline
& $\exists x~\left(  M(x)\wedge Q(x)\right)  $%
\end{tabular}
\]
perquè es considera com a prova la deducció següent:%
\[%
\begin{tabular}
[c]{lll}%
1. & $\exists x~\left(  P(x)\wedge Q(x)\right)  $ & P 1\\
2. & $\exists x$ $M(x)$ & P 2\\
3. & $P(a)\wedge Q(a)$ & EQE 1\\
4. & $Q(a)$ & EC 3\\
5. & $M(a)\wedge\lnot P(a)$ & EQE 2\\
6. & $M(a)$ & EC 5\\
7. & $M(a)\wedge Q(a)$ & IC (4,5)\\\hline
8. & $\exists x~\left(  M(x)\wedge Q(x)\right)  $ & IQE 6
\end{tabular}
\]


\begin{enumerate}
\opcio{a}{ '{E}s correcte.

}
\opcio{b}{ No és vàlid perquè EQE 2 no permet escriure $M(a)\wedge\lnot
P(a)$ sinó, per exemple, $M(b)\wedge\lnot P(b)$ essent $b$ una instància
de $x$ que no apareix abans. 

}
\opcio{c}{ El raonament no és vàlid perquè és inconsistent; de les
premisses es dedueix $P(a)\wedge\lnot P(a)$.

}
\opcio{d}{ Cap de les respostes anteriors és certa.
}
\end{enumerate}
\feedback

\pregunta Dels següents raonaments: (i) Qui menysprea a tots els fanàtics
menysprea també a tots els polítics. Algú no menysprea a un
determinat polític. Per consegüent, hi ha un fanàtic al qual no
tothom menysprea; i (2) A tots els gats simpàtics i
intel·ligents els agrada el fetge picat. Cada siamès
és simpàtic. Hi ha un siamès que no li agrada el fetge picat. Per
consegüent, hi ha un gat estúpid.

\begin{enumerate}
\opcio{a}{ (i) i (ii) són vàlids. 

}
\opcio{b}{ Cap dels dos és vàlid.

}
\opcio{c}{ (i) és vàlid, i (ii) no.

}
\opcio{d}{ (i) no és vàlid, i (ii) sí.
}
\end{enumerate}
\feedback
\end{enumerate}